\documentclass[12pt,a4paper]{article}

% ---------- Мова та шрифти ----------
\usepackage{fontspec}
\usepackage{polyglossia}
\setmainlanguage{ukrainian}
\setotherlanguage{english}

\setmainfont{DejaVu Serif}
\setsansfont{DejaVu Sans}
\setmonofont{DejaVu Sans Mono}

% ---------- Компонування ----------
\usepackage[a4paper,margin=2.2cm]{geometry}
\usepackage{graphicx}
\usepackage{float}
\usepackage{caption}
\usepackage{xcolor}
\usepackage{hyperref}
\usepackage{enumitem}
\usepackage{titlesec}
\usepackage{longtable}
\usepackage{array}
\usepackage{pdflscape}

\hypersetup{
    colorlinks=true,
    linkcolor=blue!50!black,
    urlcolor=blue!50!black,
    citecolor=blue!50!black,
    pdftitle={UML-діаграми проєкту Remote File Folder Client},
    pdfauthor={Yana Utochkina}
}

\titleformat{\section}{\normalfont\Large\bfseries}{\thesection.}{0.6em}{}
\titleformat{\subsection}{\normalfont\large\bfseries}{\thesubsection}{0.6em}{}

\setlength{\parindent}{0pt}
\setlength{\parskip}{0.6em}

\newcommand{\diagram}[4]{%
    \begin{figure}[H]
        \centering
        \includegraphics[width=#3\linewidth,height=0.8\textheight,keepaspectratio]{#1}
        \caption{#2}
        \label{#4}
    \end{figure}
}

% Варіант для широких діаграм (класи, sequence): розміщує рисунок на
% альбомній сторінці, щоб мати значно більше горизонтального простору
% і уникнути надто дрібного масштабування вмісту.
\newcommand{\diagramland}[4]{%
    \begin{landscape}
        \begin{figure}[H]
            \centering
            \includegraphics[width=#3\linewidth,height=0.85\textheight,keepaspectratio]{#1}
            \caption{#2}
            \label{#4}
        \end{figure}
    \end{landscape}
}

\title{\textbf{UML-діаграми проєкту\\[2pt] \large Remote File Folder Client\\(спрощений аналог Google Drive)}}
\author{Уточкіна Яна}
\date{Етап: Проєктування архітектури та поведінки системи \\ \today}

\begin{document}

\maketitle

\begin{abstract}
\noindent
У документі представлено комплект UML-діаграм для проєкту \textbf{Remote File Folder Client} --- застосунку для роботи з віддаленою файловою текою (спрощений аналог Google Drive) з веб- та десктоп-клієнтами, Go-бекендом, PostgreSQL та синхронізацією файлів за принципом diff-маніфестів. Діаграми впорядковано відповідно до життєвого циклу проєктування: від аналізу вимог і високорівневої архітектури до деталізованого проєктування окремих підсистем. Усі діаграми побудовано в нотації Mermaid.
\end{abstract}

\tableofcontents
\newpage

% =====================================================================
\section{Аналіз вимог}
% =====================================================================

Система надає користувачу автентифікований доступ до персональної віддаленої теки (``віртуального диска''), у якій він може переглядати список файлів разом з їх атрибутами (ім'я, дата створення, дата модифікації, автор вивантаження, автор останньої редакції), вивантажувати, завантажувати та видаляти файли, а також синхронізувати обраний локальний каталог із віддаленим.

Ключові функціональні вимоги:

\begin{itemize}[nosep]
    \item автентифікація користувача на віддаленому сервері;
    \item перегляд атрибутів файлів із можливістю показу/приховування стовпців (стовпець імені файлу приховати не можна);
    \item перегляд вмісту файлів окремих типів: \texttt{.java} -- як текст, \texttt{.png} -- як зображення;
    \item сортування файлів за іменем користувача, який востаннє редагував файл (за зростанням / спаданням);
    \item фільтрація файлів: усі файли / лише \texttt{.cs} / лише \texttt{.jpg};
    \item синхронізація локальної теки з віддаленою на основі порівняння маніфестів (\texttt{name, size, mtime, sha256}) та обробка конфліктів редагування.
\end{itemize}

Архітектурно система складається з React-клієнта (веб), Swift-клієнта (macOS, з повноцінною фоновою синхронізацією через \texttt{FSEvents}), Go API-сервера (шари handler / service / repository), PostgreSQL для метаданих та абстрактного сховища файлів (локальний диск або S3-сумісне сховище), розгорнутих на домашньому Ubuntu-сервері та опублікованих через Cloudflare Tunnel.

На основі цих вимог у наступних розділах побудовано діаграму прецедентів, а також архітектурні та деталізовані діаграми, що описують структуру і поведінку системи.

% =====================================================================
\section{Архітектурне проєктування (високий рівень)}
% =====================================================================

На цьому етапі визначається загальна архітектура системи: хто є її користувачами і які сценарії вона підтримує (діаграма прецедентів), яка загальна структура даних (діаграма класів високого рівня) та яка фізична топологія розгортання (діаграма розгортання).

% ---------------------------------------------------------------------
\subsection{Діаграма прецедентів}
% ---------------------------------------------------------------------

Діаграма (рис.~\ref{fig:usecase}) показує єдиного актора \emph{Користувача}, який працює із системою через веб- або десктоп-клієнт, та основні прецеденти, що відповідають функціональним вимогам розділу~1. Частина операцій над файлами \emph{включають} (\texttt{«include»}) прецедент \emph{``Автентифікуватися''}, оскільки доступ до робочої області можливий лише після успішного входу. Перегляд вмісту файлу \emph{розширюється} (\texttt{«extend»}) двома альтернативними сценаріями відображення (\texttt{.java} / \texttt{.png}), а синхронізація -- сценарієм розв'язання конфлікту, що активується лише за наявності розбіжностей.

\emph{Примітка щодо нотації:} Mermaid не має вбудованого типу діаграми прецедентів, тому актора, прецеденти та межу системи побудовано засобами \texttt{flowchart} (актор -- прямокутник зі заокругленими кутами, прецеденти -- ``stadium''-форма, межа системи -- \texttt{subgraph}), із явним підписом стереотипів \texttt{«include»}/\texttt{«extend»} на зв'язках.

\diagram{img/ua/01-usecase.pdf}{Діаграма прецедентів системи Remote File Folder Client}{0.95}{fig:usecase}

% ---------------------------------------------------------------------
\subsection{Діаграма класів (загальна структура)}
% ---------------------------------------------------------------------

Рис.~\ref{fig:class-hl} відображає загальну структуру даних та основні взаємозв'язки між доменними класами системи, без деталізації методів окремих сценаріїв: сутності \texttt{User} та \texttt{File}, сервісний шар (\texttt{AuthService}, \texttt{FileService}, \texttt{SyncDiffService}), шар доступу до даних (\texttt{FileRepository}) та абстракцію сховища файлів \texttt{FileStorage}, реалізовану локальним сховищем на диску та S3-сумісним сховищем -- відповідно до принципу ``сховище за інтерфейсом'', закладеного в вимогах (розділ 5.2, 5.5).

\diagramland{img/ua/02-class-highlevel.pdf}{Діаграма класів високого рівня (доменна модель)}{0.98}{fig:class-hl}

% ---------------------------------------------------------------------
\subsection{Діаграма розгортання (загальна архітектура)}
% ---------------------------------------------------------------------

Рис.~\ref{fig:deployment} відображає фізичну топологію системи: клієнтські пристрої (браузер, macOS), периферію Cloudflare (термінація TLS, публічний DNS) та домашній Ubuntu-сервер, що через Docker Compose підіймає три контейнери -- \texttt{cloudflared} (вихідний тунель без відкриття портів), \texttt{api} (Go-бінарник) та \texttt{db} (PostgreSQL) -- разом із постійним томом \texttt{pgdata} та bind-mount каталогом для файлового сховища.

\diagram{img/ua/10-deployment.pdf}{Діаграма розгортання}{0.98}{fig:deployment}

% =====================================================================
\section{Деталізоване проєктування}
% =====================================================================

Деталізоване проєктування розкриває внутрішню структуру та поведінку окремих підсистем: класи, що беруть участь у реалізації конкретного прецедента (VOPC-діаграма), конкретний знімок станів об'єктів під час конфлікту (діаграма об'єктів), взаємодію об'єктів у часі -- як у часовій (sequence), так і в структурній (communication) нотації -- поведінку файлу залежно від стану синхронізації (діаграма станів), структуру програмних компонентів (діаграма компонентів), організацію коду за модулями (діаграма пакетів) та логіку бізнес-процесу синхронізації (діаграма активності).

% ---------------------------------------------------------------------
\subsection{VOPC-діаграма класів для прецедента ``Синхронізувати локальну теку з віддаленою''}
% ---------------------------------------------------------------------

Згідно з варіантом завдання, як прецедент для VOPC-діаграми (\emph{View Of Participating Classes}) обрано \textbf{``Синхронізувати локальну теку з віддаленою''} -- найскладніший і найбільш детально описаний у вимогах сценарій (розділ 6 вимог). Класи розподілено за стандартними стереотипами аналізу:

\begin{itemize}[nosep]
    \item \texttt{«boundary»} -- елементи взаємодії з користувачем: спостерігач за текою (\texttt{FolderWatcherUI}), кнопка ручної синхронізації (\texttt{SyncNowButtonUI}), діалог розв'язання конфлікту (\texttt{ConflictDialogUI});
    \item \texttt{«control»} -- керуюча логіка: \texttt{SyncController} (клієнт), \texttt{SyncDiffService} (сервер, обчислення diff), \texttt{ConflictResolver} (застосування вибору користувача);
    \item \texttt{«entity»} -- дані предметної області: \texttt{ManifestEntry}, \texttt{SyncCursor}, \texttt{File}, \texttt{ChangeLogEntry}.
\end{itemize}

\diagramland{img/ua/03-class-vopc-sync.pdf}{VOPC-діаграма класів для прецедента ``Синхронізувати локальну теку з віддаленою''}{0.98}{fig:vopc}

% ---------------------------------------------------------------------
\subsection{Діаграма об'єктів (знімок конфлікту)}
% ---------------------------------------------------------------------

На відміну від класових діаграм, що описують структуру типів, діаграма об'єктів (рис.~\ref{fig:object}) фіксує \emph{конкретний знімок} екземплярів під час виконання -- одну мить сценарію синхронізації, коли виявлено конфлікт (розділ 6.2 вимог, випадок ``exists, hash differs''). Локальний екземпляр \texttt{localFile} (версія 44, обчислена на клієнті) та серверний \texttt{remoteFile} (версія 45) мають однакове ім'я й розмір-ідентичність за \texttt{id}, але різняться за \texttt{sha256} -- це і є ознака редагування з обох боків. Запис \texttt{cl1 : ChangeLogEntry} підтверджує, що версія 45 новіша за \texttt{lastSyncedVersion = 42}, збережену в курсорі клієнта, тобто зміна на сервері відбулася вже після останньої успішної синхронізації.

\emph{Примітка щодо нотації:} Mermaid не має вбудованого типу діаграми об'єктів, тому екземпляри побудовано засобами \texttt{flowchart}: заголовок кожного об'єкта підписано за конвенцією UML як \underline{ім'я\_екземпляра~:~Клас}, а нижче через розділювач перелічено конкретні значення атрибутів.

\diagramland{img/ua/12-object-conflict.pdf}{Діаграма об'єктів: знімок стану під час конфлікту редагування}{0.9}{fig:object}

% ---------------------------------------------------------------------
\subsection{Діаграми взаємодії (Sequence \& Communication Diagrams)}
% ---------------------------------------------------------------------

Наведено три діаграми послідовності, що ілюструють ключові сценарії взаємодії клієнта, Go API-сервера та бази даних, а також комунікаційну діаграму як альтернативне представлення одного з цих сценаріїв.

\subsubsection*{Автентифікація}

Рис.~\ref{fig:seq-login} показує послідовність перевірки облікових даних користувача та видачі пари токенів (access/refresh) відповідно до опису автентифікації у вимогах (розділ 1, 5.2).

\diagramland{img/ua/05-seq-login.pdf}{Діаграма послідовності: автентифікація користувача}{0.95}{fig:seq-login}

\subsubsection*{Вивантаження файлу}

Рис.~\ref{fig:seq-upload} ілюструє шлях запиту від клієнта через шар обробників (\texttt{handler}) до сервісного шару, який послідовно записує вміст файлу у сховище (\texttt{FileStorage}), а після підтвердження запису -- метадані файлу у PostgreSQL через \texttt{FileRepository}.

\diagramland{img/ua/06-seq-upload.pdf}{Діаграма послідовності: вивантаження файлу на сервер}{0.95}{fig:seq-upload}

\subsubsection*{Синхронізація та розв'язання конфлікту}

Рис.~\ref{fig:seq-sync} деталізує протокол \texttt{POST /api/sync/diff} (розділ 6.6 вимог): клієнт надсилає локальний маніфест і \texttt{lastSyncedVersion}, сервер повертає перелік дій для кожного файлу, а конфліктні файли обробляються через запит вибору користувача (keep local / keep remote / keep both).

\diagramland{img/ua/07-seq-sync.pdf}{Діаграма послідовності: синхронізація та розв'язання конфлікту}{0.98}{fig:seq-sync}

\subsubsection*{Той самий сценарій у нотації комунікаційної діаграми}

Рис.~\ref{fig:comm-sync} зображує \textbf{той самий} сценарій синхронізації та розв'язання конфлікту, що й рис.~\ref{fig:seq-sync}, але в нотації комунікаційної (communication) діаграми: акцент зміщено з часової осі на статичні зв'язки між об'єктами, а порядок повідомлень передається нумерацією (з десятковою вкладеністю для повідомлень усередині умовних гілок -- наприклад, \texttt{5.1--5.3} для взаємодії з користувачем у разі конфлікту). Обидві діаграми узгоджені між собою -- використовують ті самі об'єкти й повідомлення -- і демонструють два стандартні способи показати одну взаємодію.

\emph{Примітка щодо нотації:} Mermaid не має вбудованого типу комунікаційної діаграми, тому вона побудована засобами \texttt{flowchart}: об'єкти -- вузли графа, а нумеровані підписи на зв'язках відповідають повідомленням.

\diagramland{img/ua/11-comm-sync.pdf}{Комунікаційна діаграма: синхронізація та розв'язання конфлікту}{0.95}{fig:comm-sync}

% ---------------------------------------------------------------------
\subsection{Діаграма станів}
% ---------------------------------------------------------------------

Рис.~\ref{fig:state} описує поведінку сутності \emph{File} залежно від стану синхронізації: від локального/віддаленого створення, через передачу (\texttt{Uploading}/\texttt{Downloading}), до синхронізованого стану, можливих локальних чи віддалених змін і, за одночасної зміни з обох боків, -- конфлікту, що розв'язується користувачем. Діаграма безпосередньо реалізує чотири випадки порівняння маніфестів із розділу 6.2 вимог (upload / download / no-op / conflict) разом з обробкою видалень через версійний курсор (розділ 6.3).

\diagramland{img/ua/08-state-file.pdf}{Діаграма станів файлу в процесі синхронізації}{0.85}{fig:state}

% ---------------------------------------------------------------------
\subsection{Діаграма компонентів}
% ---------------------------------------------------------------------

Рис.~\ref{fig:component} деталізує програмні компоненти системи та їх залежності: клієнтські застосунки, тунель \texttt{cloudflared}, а всередині Go API-сервера -- шари \texttt{Handler} $\to$ \texttt{Service} $\to$ \texttt{Repository}, а також незалежний від конкретної реалізації інтерфейс \texttt{FileStorage}, що дозволяє замінити локальне сховище на S3-сумісне без зміни бізнес-логіки (розділ 5.2, 5.5 вимог).

\diagram{img/ua/09-component.pdf}{Діаграма компонентів системи}{0.9}{fig:component}

% ---------------------------------------------------------------------
\subsection{Діаграма пакетів}
% ---------------------------------------------------------------------

Якщо діаграма компонентів (рис.~\ref{fig:component}) показує компоненти часу виконання та інтерфейси між ними, то діаграма пакетів (рис.~\ref{fig:package}) відображає \emph{статичну} організацію коду за модулями відповідно до запланованої структури репозиторію (\texttt{backend}, \texttt{frontend}, \texttt{desktop}, \texttt{migrations}): усередині Go-модуля -- шарова структура пакетів \texttt{handler} $\to$ \texttt{service} $\to$ \texttt{repo}/\texttt{storage}, що відповідає конвенції ``жодних прямих викликів файлової системи з handler'' та принципу доступу до сховища лише через інтерфейс \texttt{FileStorage}. Пакет \texttt{repo} додатково залежить від схеми, визначеної в \texttt{migrations}.

\emph{Примітка щодо нотації:} Mermaid не має вбудованого типу діаграми пакетів, тому пакети зображено вкладеними \texttt{subgraph}-блоками (стереотип \texttt{«package»} у заголовку), а залежності -- пунктирними стрілками з підписом дії.

\diagram{img/ua/13-package-arch.pdf}{Діаграма пакетів: організація коду за модулями}{0.95}{fig:package}

% ---------------------------------------------------------------------
\subsection{Діаграма активності}
% ---------------------------------------------------------------------

Рис.~\ref{fig:activity} деталізує бізнес-логіку одного циклу синхронізації: обчислення локального маніфесту, надсилання його на сервер, визначення дії для кожного файлу (upload/download/no-op/conflict/delete\_local), виконання дії або, у разі конфлікту, отримання рішення від користувача (keep local / keep remote / keep both) і, зрештою, збереження нової версії курсора синхронізації (розділ 6 вимог).

\diagram{img/ua/04-activity-sync.pdf}{Діаграма активності: цикл синхронізації}{0.98}{fig:activity}

% =====================================================================
\section{Висновки}
% =====================================================================

Розроблений комплект діаграм послідовно розкриває систему Remote File Folder Client від бізнес-вимог до деталізованої структури і поведінки: діаграма прецедентів фіксує функціональні сценарії користувача; загальна діаграма класів і діаграма розгортання задають архітектурний каркас; VOPC-діаграма, діаграма об'єктів, діаграми послідовності та комунікаційна, станів, компонентів, пакетів та активності деталізують реалізацію найскладнішого сценарію -- синхронізації файлів між локальною та віддаленою текою, включно з протоколом порівняння маніфестів та розв'язанням конфліктів редагування. Зокрема, діаграма об'єктів і комунікаційна діаграма доповнюють відповідно VOPC-діаграму та діаграми послідовності конкретним знімком станів та альтернативною (структурною, а не часовою) нотацією того самого сценарію, а діаграма пакетів фіксує статичну модульну організацію коду, узгоджену з плановою структурою репозиторію. Побудовані діаграми узгоджені між собою (спільні класи, сутності та дії) та слугують основою для подальшої розробки й тестування системи.

\end{document}
